% 文献库主文档(中文模板)
% 引擎必须用 XeLaTeX(中文);HTML 构建由 make4ht --engine xelatex 完成
% 新增文献:refs.bib 加一条 entry + keywords 写分类标签,push 后网页自动更新
\documentclass[11pt]{article}
\usepackage[UTF8]{ctex}
\usepackage[backend=biber, style=numeric, sorting=none, maxbibnames=99]{biblatex}
\usepackage[colorlinks=true, linkcolor=blue, urlcolor=blue, citecolor=blue]{hyperref}
\addbibresource{refs.bib}

% 关键:nocite 让未被正文引用过的条目也出现在各分类列表里
\nocite{*}

\title{Paper List}
\author{whoisKIWIIZZZ}
\date{\today}

\begin{document}
\maketitle
% \maketitle 会在 HTML 里丢失 xeCJK 中文(tex4ht 已知问题),改为手动标题页
% \begin{center}
% {\LARGE Paper Lists}\\[0.6ex]
% {\normalsize whoisKIWIIZZZ}\\[0.3ex]
% {\normalsize \today}
% \end{center}


% \section{如何维护这个列表}
% 每篇文献是一条 BibTeX 记录,都存在 \texttt{refs.bib} 里:
% \begin{itemize}
% \item \textbf{新增}:往 \texttt{refs.bib} 粘贴一条 BibTeX 条目(谷歌学术或 arXiv 都能导出);
% \item \textbf{分类}:给条目写 \texttt{keywords},可多个——同一篇论文会自动出现在所有对应分类下;
% \item \textbf{引用}:正文写 \verb|\parencite{yao2023react}|,渲染成 “[3]” 这样的编号,点击跳到文末对应条目;
% \item \textbf{发布}:\texttt{git push} 后 GitHub Actions 自动构建部署,不需要本地编译;
% 本地预览则运行 \texttt{./build.sh},打开 \texttt{docs/paperlist.html}。
% \end{itemize}

\section{World Model}
\subsection{JEPAs}
Training object is (likely) representations, and use CEM to MPC.
\begin{itemize}
    \item I-JEPA,V-JEPA 2506.09985,2512.10942,2404.08471,2301.08243,2603.14482
    \item leWM-based 2603.19312,2605.26379,2606.32026,
    \item DINO-WM 2411.04983
\end{itemize}
\subsection{latent representations}
\begin{itemize}
    \item 2603.12231, smoothness for representations
    \item 2607.27924, project DINO token
    \item  \textit{2609.29171, planning in latent}
\end{itemize}
\subsection{WAMs}
\subsection{Others}
\begin{itemize}
    \item 2603.04553,Point WAMs
    \item 2610.26844, dataset for other goals
\end{itemize}
\section{Loop}
\begin{itemize}
    \item 2602.02156,LoopViT
    \item 2512.19941, use DP to find  training-free loops in ViTs
    \item 2604.09168, refine in Diffusion
    \item 2609.40305-loopdit
\end{itemize}
\section{NCA}
\begin{itemize}
    \item 2604.24990, Summary
    \item 2609.01449, Diffusion as NCA 
    \item 2610.08518,2609.36126, 2D reasoning w/ NCA
\end{itemize}
\section{Theory}
\begin{itemize}
    \item 2202.10054,question ID/OOD acc using gradient flow. Check if u wish to see analyses of subspace
    \item 2401.17505,great work of "capacity" of models\footnote{here to say, inductive bias is needed for models to learn}
    \item several works of JEPAs.
    \begin{itemize}
        \item 2605.26379,Gaussian world, how SIGReg work ?
        \item 2608.22764, LpWM , sprase 
        \item 2603.12231, smoothness for representations
    \end{itemize}
\end{itemize}
% \section{Agent}
% ReAct 让大模型在推理中交替“隐式思考”和“外部行动” \parencite{yao2023react};
% Toolformer 则让模型自己学会在合适的位置调用工具 \parencite{schick2023toolformer}。
% \printbibliography[keyword=agent, heading=subbibliography, title={Agent 方法}]

% \section{RAG}
% 检索增强生成把外部知识库引入生成过程 \parencite{lewis2020rag}。
% \printbibliography[keyword=rag, heading=subbibliography, title={RAG 论文}]

% \section{经典基础}
% \printbibliography[keyword=classic, heading=subbibliography, title={必读经典}]

% \section{Scaling}
% Scaling 经验律:模型规模、数据量与算力的权衡 \parencite{kaplan2020scaling}。
% \printbibliography[keyword=scaling, heading=subbibliography, title={Scaling laws}]

% \section{全部文献}
% \printbibliography[heading=subbibliography, title={按年份排序}]

\end{document}
